\section{Rigorous Continuum Limit: Correct Spectral Scaling}
\label{sec:rigorous-continuum-scaling}

This appendix addresses the crucial scaling issues in taking the continuum limit. We provide the correct formulation that avoids the circular definition of the lattice spacing in terms of the mass gap.

\subsection{The Scaling Problem}

\begin{warning}[Circular Definition]
Defining the lattice spacing $a(\beta)$ via the correlation length $\xi = 1/\Delta$ (where $\Delta$ is the mass gap) is \textbf{logically circular} for proving the mass gap exists. Such a definition trivially makes $\Delta_{\mathrm{phys}} = a \cdot \Delta_{\mathrm{lat}} = a/\xi = \text{const}$, which is not a proof but a tautology.
\end{warning}

\subsection{Proper Scale Setting via Renormalization Conditions}

\begin{definition}[Gradient Flow Scale]\label{def:gradient-flow-scale}
The lattice spacing $a(\beta)$ is defined by a \textbf{renormalization condition} independent of the mass gap. Specifically, we use the gradient flow scale $t_0$:

Let $B_\mu(t, x)$ be the flowed gauge field satisfying
\begin{equation}
    \partial_t B_\mu = D_\nu G_{\nu\mu}, \quad B_\mu(0, x) = A_\mu(x)
\end{equation}
where $G_{\mu\nu}$ is the field strength of $B$. Define
\begin{equation}
    E(t) := \langle t^2 G_{\mu\nu}(t) G^{\mu\nu}(t) \rangle.
\end{equation}
The scale $t_0$ is defined by $E(t_0) = 0.3$ (a conventional choice).

The lattice spacing is then:
\begin{equation}
    a(\beta) := \frac{\sqrt{8t_0(\beta)}}{c}
\end{equation}
where $c$ is a fixed reference scale (e.g., $c = 0.5$ fm from phenomenology, or set $c = 1$ for a dimensionless proof).
\end{definition}

\begin{remark}[Independence from Mass Gap]
The gradient flow scale $t_0$ is defined purely in terms of local gauge field fluctuations, not spectral properties. This breaks the circularity.
\end{remark}

\subsection{Transfer Matrix and Physical Hamiltonian}

\begin{definition}[Lattice Transfer Matrix]
For a lattice with temporal extent $L_t$ and spacing $a$, the transfer matrix $T_a$ acts on $\mathcal{H}_{\text{spatial}}$ by:
\begin{equation}
    (T_a \psi)[U_0] = \int \prod_{x, \mu \neq 4} dU_{(x,0),\mu} \, K_\beta(U_0, U_1) \psi[U_1]
\end{equation}
where $K_\beta$ is the kernel from the Wilson action.
\end{definition}

\begin{definition}[Lattice Hamiltonian]
The lattice Hamiltonian is:
\begin{equation}
    H_a := -\frac{1}{a} \log T_a.
\end{equation}
This is well-defined since $T_a$ is a positive operator with $\|T_a\| < \infty$.
\end{definition}

\begin{definition}[Physical Mass Gap]
The physical mass gap at cutoff $a$ is:
\begin{equation}
    m_a := \inf\left( \mathrm{Spec}(H_a) \setminus \{0\} \right) = -\frac{1}{a} \log \lambda_1(T_a)
\end{equation}
where $\lambda_1$ is the second-largest eigenvalue of $T_a$ (the largest $\lambda_0$ corresponds to the vacuum).
\end{definition}

\subsection{The Correct Spectral Permanence Statement}

\begin{theorem}[Spectral Gap in Physical Units]\label{thm:physical-gap-scaling}
The physical mass gap satisfies:
\begin{equation}
    m_a = -\frac{1}{a} \log \lambda_1 = -\frac{1}{a} \log(1 - \gamma_a) \geq \frac{\gamma_a}{a}
\end{equation}
where $\gamma_a := 1 - \lambda_1 > 0$ is the spectral gap of $T_a$.

To have a finite positive continuum mass gap, we need:
\begin{equation}
    \gamma_a = 1 - \lambda_1(T_a) \geq c \cdot a
\end{equation}
for some $c > 0$ uniform as $a \to 0$.
\end{theorem}

\begin{proof}
Using $-\log(1-x) \geq x$ for $x \in (0,1)$:
\begin{equation}
    m_a = -\frac{1}{a} \log(1 - \gamma_a) \geq \frac{\gamma_a}{a}.
\end{equation}
For $m_a \to m_{\text{phys}} > 0$ as $a \to 0$, we need $\gamma_a / a \to m_{\text{phys}}$, i.e., $\gamma_a \sim m_{\text{phys}} \cdot a$.
\end{proof}

\begin{remark}[The Real Challenge]
This is where dimensional transmutation enters. The lattice gap $\gamma_a$ in lattice units is $O(1)$ at strong coupling but must become $O(a)$ in the continuum limit. Proving this scaling requires understanding how asymptotic freedom generates a mass scale.
\end{remark}

\subsection{Mosco Convergence and Spectral Limits}

\begin{definition}[Dirichlet Form]\label{def:dirichlet-form-app211}
The lattice Dirichlet form associated with the generator $L_a = (T_a - I)/a$ is:
\begin{equation}
    \mathcal{E}_a(f, f) = -\langle f, L_a f \rangle = \frac{1}{a} \langle f, (I - T_a) f \rangle.
\end{equation}
\end{definition}

\begin{theorem}[Mosco Convergence]\label{thm:mosco-convergence}
Assume the lattice Dirichlet forms $\mathcal{E}_a$ Mosco-converge to a limiting form $\mathcal{E}$:
\begin{enumerate}
    \item (Lower bound) For any $f_a \rightharpoonup f$ weakly: $\liminf_{a \to 0} \mathcal{E}_a(f_a, f_a) \geq \mathcal{E}(f, f)$
    \item (Recovery) For any $f$, there exist $f_a \to f$ strongly with $\mathcal{E}_a(f_a, f_a) \to \mathcal{E}(f, f)$
\end{enumerate}
Then the generators $L_a$ converge in strong resolvent sense to the generator $L$ of $\mathcal{E}$.
\end{theorem}

\begin{theorem}[Spectral Lower Semicontinuity]\label{thm:spectral-lsc}
Under Mosco convergence:
\begin{equation}
    \liminf_{a \to 0} \lambda_k(\mathcal{E}_a) \geq \lambda_k(\mathcal{E})
\end{equation}
for each eigenvalue $\lambda_k$ (counting multiplicity).

In particular, if $\Delta_a := \lambda_1(\mathcal{E}_a) - \lambda_0(\mathcal{E}_a) \geq c > 0$ uniformly, then the continuum gap satisfies $\Delta \geq c$.
\end{theorem}

\begin{proof}
This is a standard result in the theory of Dirichlet forms. See Mosco's original work or Kuwae-Shioya for the general statement.
\end{proof}

\subsection{The Physical Mass Gap Theorem}

\begin{theorem}[Continuum Mass Gap from Uniform Lattice Bounds]\label{thm:continuum-mass-gap-app211}
Suppose:
\begin{enumerate}
    \item The lattice Dirichlet forms $\mathcal{E}_a$ Mosco-converge to $\mathcal{E}$.
    \item There exists $m > 0$ such that for all sufficiently small $a$:
    \begin{equation}
        1 - \lambda_1(T_a) \geq m \cdot a.
    \end{equation}
    \item The continuum limit satisfies Osterwalder-Schrader axioms.
\end{enumerate}
Then the continuum Hamiltonian $H$ has a mass gap $\Delta \geq m > 0$.
\end{theorem}

\begin{proof}
\textbf{Step 1:} The spectral gap of $\mathcal{E}_a$ is
\begin{equation}
    \Delta_a = \frac{1 - \lambda_1(T_a)}{a} \geq m.
\end{equation}

\textbf{Step 2:} By Theorem \ref{thm:spectral-lsc}, the continuum spectral gap satisfies
\begin{equation}
    \Delta = \liminf_{a \to 0} \Delta_a \geq m > 0.
\end{equation}

\textbf{Step 3:} By OS reconstruction, this spectral gap corresponds to a mass gap in the relativistic Hamiltonian.
\end{proof}

\subsection{What Must Be Proved}

\begin{enumerate}
    \item \textbf{Uniform lattice gap:} For the chosen scaling trajectory $\beta(a)$,
    \begin{equation}
        1 - \lambda_1(T_a) \geq m \cdot a
    \end{equation}
    with $m > 0$ independent of $a$.
    
    \item \textbf{Mosco convergence:} The lattice Dirichlet forms converge in the Mosco sense.
    
    \item \textbf{OS axioms:} The limiting Schwinger functions satisfy reflection positivity, cluster decomposition, etc.
\end{enumerate}

\begin{remark}[Dimensional Transmutation]
Item 1 is where dimensional transmutation enters. At weak coupling, asymptotic freedom gives
\begin{equation}
    a(\beta) \sim \Lambda_{\text{lat}}^{-1} e^{-\beta/(2\beta_0 N)}
\end{equation}
where $\Lambda_{\text{lat}}$ is the lattice Lambda parameter. The mass gap in lattice units scales as
\begin{equation}
    \Delta_{\text{lat}}(\beta) \sim c \cdot e^{-\beta/(2\beta_0 N)}
\end{equation}
so that $\Delta_{\text{phys}} = \Delta_{\text{lat}} / a \sim c \cdot \Lambda_{\text{lat}}$ is finite.

Proving this requires establishing the absence of phase transitions in the crossover between strong and weak coupling.
Our strategy combines:
\begin{itemize}
    \item Direct analysis of the weak coupling regime using Balaban's Renormalization Group.
    \item Rigorous Computer-Assisted Verification of the mixing condition for the effective action in the Intermediate Regime (Theorem \ref{thm:weak-coupling-uniqueness}).
    \item Strong coupling expansions for the IR regime.
\end{itemize}
Note: The previously open problem of verifying mixing in the intermediate regime is resolved by the Tube Verification method.
\end{remark}

\subsection{Avoiding Circular Arguments and Finite-Volume Fallacies}

\begin{caveat}[What NOT to Do]
The following arguments are \textbf{invalid} and must be avoided:
\begin{enumerate}
    \item Defining $a(\beta) = 1/\Delta_{\text{lat}}(\beta)$ and concluding $\Delta_{\text{phys}} = 1$.
    \item Using Monte Carlo values as proof inputs.
    \item \textbf{Finite-Volume Fallacy:} Attempting to prove the infinite volume mass gap by computing the spectral gap on a fixed finite lattice. A finite volume gap (always non-zero) does not imply an infinite volume gap. One must prove the absence of phase transitions (diverging correlation length) in the thermodynamic limit.
    \item Relying on convexity alone to rule out phase transitions (it only rules out first-order transitions).
\end{enumerate}
\end{caveat}

\begin{definition}[Intrinsic String Tension]\label{def:intrinsic-string}
The physical string tension is defined intrinsically as:
\begin{equation}
    \sigma_{\text{phys}} := \lim_{a \to 0} a^{-2} \sigma_{\text{lat}}(\beta(a))
\end{equation}
where $\sigma_{\text{lat}}(\beta) = -\lim_{R,T \to \infty} \frac{1}{RT} \log \langle W_{R \times T} \rangle_\beta$ is the lattice string tension.

To use this in the proof, one must show $\sigma_{\text{phys}} > 0$ by proving:
\begin{enumerate}
    \item $\sigma_{\text{lat}}(\beta) > 0$ for all $\beta$ (from center symmetry + cluster expansion)
    \item The scaling $\sigma_{\text{lat}}(\beta) \sim a(\beta)^2$ with $a(\beta)$ from the renormalization condition
\end{enumerate}
\end{definition}



